Adversarial work removes the assumption that unreliable evidence arrives by accident and
replaces it with an attacker who chooses the evidence (Appendix~\ref{app:adversarial}).
Because RAG externalises knowledge into corpora, indices, prompt templates and tools, every
one of these becomes an attack surface that a purely parametric model does not have.

\paragraph{Poisoning turns small corpus edits into large output changes.}
\citet{poisoningRetrievalCorpora2023} showed that a handful of crafted passages is enough to
distort a dense retriever's neighbourhood structure, and PoisonedRAG carried the same idea
through to the generated answer \citep{poisonedRAG2024}. Later work varies the trigger and the
persistence model: BadRAG studies triggered retrieval \citep{badRAG2024}, while TrojanRAG and
Phantom induce attacker-chosen behaviour through retrieved context alone, without touching
model weights \citep{trojanRAG2024,phantom2024}. What these share is an amplification effect:
retrieval converts a marginal perturbation of the corpus into a first-class input to the model.

\paragraph{Injection and poisoning are different threats.}
Indirect prompt injection plants \emph{executable instructions} rather than false evidence, and
the payload must first clear a retrieval barrier before it can act \citep{rrag005}; AIP instead
exploits prompts shared across users \citep{aip2025}. The distinction matters for defence
design, because a filter tuned to detect implausible claims will not detect a plausible
document that ends with an instruction. Treating templates, ingestion pipelines and retrieved
text as untrusted boundaries is therefore a stronger design stance than treating the corpus as
the only thing to sanitise.

\paragraph{Defences are not comparable across threat models.}
Proposed defences span certified aggregation \citep{certifiedRAG2024}, adversarial reranking
\citep{rrag004}, attention-level filtering \citep{rrag015}, and post-hoc traceback of the
responsible passage \citep{rrag128}, and each is evaluated against the threat its authors
assumed. Two further sources of divergence are less visible. Papers differ in what the
attacker is allowed to know and do, and they differ in the unit that the attack budget is
counted in---documents injected into the corpus versus documents actually retrieved into the
context---so nominally identical budgets can differ by an order of magnitude. Reporting
adversary access and budget scope, coordinate C2 in \S\ref{sec:coordinates}, is what makes two
attack success rates commensurable. Recent evaluation suites move this way, fixing
corpus, retriever and injection budget before pitting attacks against one another
\citep{rrag019,rrag023,rrag033}.
